\documentclass[12pt,a4paper,oneside]{article}
\usepackage[brazil]{babel}
\usepackage[utf8]{inputenc}
\usepackage{olymp}

\setcounter{problem}{0}

\begin{document}

\begin{problem}{Candy Crush}{frutas}{2}

Candy Crush é um jogo popular em que se deve trocar doces vizinhos para formar fileiras de três ou mais doces iguais.
Assim, quando há dois doces iguais um ao lado do outro, é interessante colocar mais um igual na fileira para eliminá-los.

Nesse exercício você deve apenas identificar e contar quantos doces têm um vizinho igual,
considerando como vizinhos os que estão imediatamente acima, abaixo, à esquerda ou à direita.

A figura abaixo mostra uma versão com frutas. As 10 frutas destacadas são as que têm vizinha igual.

\begin{center}
\hfill
\includegraphics[scale=0.325]{figuras/frutas.png}
\hfill
\includegraphics[scale=0.325]{figuras/frutas2.png}
\hfill{ }
\end{center}

%\Notes
%
%Observações, se tiver.

\InputFile

A entrada começa com dois inteiros, $M$ e $N$, indicando respectivamente o número de linhas e colunas do tabuleiro.
Em seguida virão $M$ linhas, cada uma com $N$ números inteiros entre 0 e 9, representando as diferentes frutas do jogo.
Restrições: $1 \leq M, N \leq 100$.

\OutputFile

Seu programa deve escrever o número de frutas com pelo menos uma vizinha igual.

%\Notes

\Example

\begin{example}
\exmp{
5 5
4 2 2 2 3
0 9 7 6 5
0 4 4 7 8
3 5 7 1 1
7 3 8 1 5
}{
10
}%
\end{example}

\begin{example}
\exmp{
4 4
1 1 1 2
2 3 1 1
2 1 2 2
1 1 1 1
}{
14
}%
\end{example}

%Comentários sobre o exemplo, se necessário.
O primeiro exemplo corresponde ao da imagem mostrada no enunciado (maçã = 4, banana = 2, ...).

\end{problem}

\end{document}